\section{Analysis Requirements}\label{sec:Requirements}

Designing a static analysis to detect missing \texttt{RB\_GC\_GUARD} directives requires addressing several fundamental requirements. This section characterizes what any such analysis must accomplish, informing both our choice of analysis framework and the design of our detection algorithm.

\subsection{Overview of Requirements}

To detect omitted guards, an analyzer must satisfy five requirements:

\begin{description}
\item[R1: GC-Triggering Function Identification.] The analyzer must identify all functions that may invoke the garbage collector, either directly or through transitive calls. This requires computing reachability in the call graph from every function to CRuby's GC entry point.

\item[R2: Subordinate Pointer Derivation Detection.] The analyzer must recognize program points where a raw C pointer is derived from a Ruby \texttt{VALUE}. These derivation sites include macro invocations such as \texttt{RSTRING\_PTR} and \texttt{TypedData\_Get\_Struct}, as well as function calls that return pointers into object-owned memory.

\item[R3: Ownership Binding Between Pointers and VALUEs.] The analyzer must establish and track the ownership relationship between each subordinate pointer and its owning \texttt{VALUE} variable. When a derivation occurs, a logical binding is created: the validity of the subordinate pointer depends on the liveness of the owning \texttt{VALUE}.

\item[R4: Control-Flow Analysis of Accesses.] The analyzer must determine the temporal ordering among three kinds of program points: accesses to subordinate pointers, accesses to owning \texttt{VALUE} variables, and GC-triggering calls. A hazard exists when a subordinate pointer is used after a GC trigger, while the owning \texttt{VALUE} has no access after the trigger.

\item[R5: Guard Existence Checking.] The analyzer must determine whether the owning \texttt{VALUE} is protected by an explicit \texttt{RB\_GC\_GUARD} directive. Protection can also come from a subsequent \texttt{VALUE}-level access that forces the compiler to keep it live. When neither protection exists, a missing guard should be flagged.
\end{description}

Combining these requirements, a missing guard is indicated when: a subordinate pointer $p$ is obtained from a \texttt{VALUE} $v$ (R2, R3), a GC-triggering call occurs after the derivation (R1), the pointer $p$ is used after the trigger (R4), and the variable $v$ has no \texttt{VALUE}-level access after the trigger and no explicit guard (R4, R5).

\subsection{Challenges in Meeting the Requirements}

Among these five requirements, three present significant technical challenges that influence tool selection.

\subsubsection*{Challenge of R1: Recursive Call-Graph Reachability.}
Identifying GC-triggering functions requires computing the transitive closure of the call graph rooted at CRuby's GC entry point \texttt{gc\_enter}. The Ruby runtime is deeply interconnected, and many non-trivial operations can transitively reach the GC. A simple heuristic (e.g., ``only \texttt{rb\_gc} triggers GC'') would miss the vast majority of real hazards, since object allocation, method dispatch, and string operations may all trigger GC when memory pressure is high. Computing this transitive closure efficiently over a large call graph requires recursive query evaluation with proper memoization.

\subsubsection*{Challenge of R3: Tracking Ownership Bindings.}
Subordinate pointers are obtained through many different APIs, each with different signatures and expansion patterns (Table~\ref{tab:derivation_patterns}). The analyzer must not only recognize these patterns but also maintain the logical binding between each derived pointer and its owning \texttt{VALUE} throughout the analysis. This binding must persist across the control-flow analysis in R4, requiring the analyzer to correlate information from different program points.

\begin{table}[ht]
\centering
\caption{Categories of derivation patterns in CRuby}
\label{tab:derivation_patterns}
\begin{tabular}{lll}
\toprule
\textbf{Category} & \textbf{Examples} & \textbf{Return Type} \\
\midrule
String access & \texttt{RSTRING\_PTR}, \texttt{StringValueCStr} & \texttt{char *} \\
Array access & \texttt{RARRAY\_PTR}, \texttt{RARRAY\_CONST\_PTR} & \texttt{VALUE *} \\
Typed data & \texttt{TypedData\_Get\_Struct}, \texttt{DATA\_PTR} & \texttt{void *} \\
Binary data & \texttt{rb\_io\_buffer\_get\_bytes} & \texttt{void *} \\
\bottomrule
\end{tabular}
\end{table}

\subsubsection*{Challenge of R4: Multi-Entity Control-Flow Ordering.}
The core detection logic requires reasoning about the temporal ordering of three distinct entities: subordinate pointer uses, owning \texttt{VALUE} accesses, and GC triggers. The analyzer must determine whether a path exists where the last \texttt{VALUE}-level access precedes a GC trigger, which in turn precedes a subordinate pointer use. This ordering constraint must be evaluated efficiently across all relevant combinations in the program.

\subsection{Scale of the Codebase}

CRuby comprises over 300,000 lines of C code with more than 23,000 \texttt{VALUE} variables and over 42,000 functions. Naive enumeration of all possible derivation-trigger-use triples would be computationally impractical: a cubic algorithm over these quantities could require examining on the order of $10^{13}$ combinations. Any practical analyzer must employ efficient indexing and pruning strategies to focus on relevant program fragments while completing in a reasonable time frame.

\subsection{Detecting Existing Guards}

The \texttt{RB\_GC\_GUARD} mechanism is implemented as a macro that expands to a GCC statement expression in preprocessed code. Listing~\ref{lst:guard_expansions} shows the expansion pattern that our analyzer must recognize.

\begin{figure}[tb]
\begin{lstlisting}[caption={Expansion of \texttt{RB\_GC\_GUARD} with GCC.},label={lst:guard_expansions}]
// Source code:
RB_GC_GUARD(str);

// Expansion:
(*__extension__ ({
    volatile VALUE *rb_gc_guarded_ptr = &(str);
    __asm__("" : : "m"(rb_gc_guarded_ptr));
    rb_gc_guarded_ptr;
}));
\end{lstlisting}
\end{figure}

Reliably detecting whether a guard is present requires recognition of this characteristic expansion pattern in the preprocessed code, specifically the assignment to a \texttt{volatile} pointer named \texttt{rb\_gc\_guarded\_ptr}.
